% 第4章 AI 预测模型、超参优化与交叉验证（占位，待 A02/A03 完成后填充）
% 素材：tcad-research-training/predictor/README.md（predict-r / predict-f 合同）；
%       experiments/dram-see-v3/r50-main/data/delivery-r4/20261012T0033Z-full/prediction-scores.csv（4 岭回归基线评分，域内 ΔM MAE≈0.53–0.55 mV）；
%       shared/metrics-and-labels.md §5（泄漏控制：同事件/配对 bit/细化归同一组，旧留出再调参不算盲测）；
%       r0.6 决定：R→A→F 残差模型，比较纯机理/浅层树/纯 MLP/RC+残差/RC+残差+KCL
% 缺口：A02 物理约束消融、A03 分组 CV 与冻结盲测
\section{AI 预测模型、超参数优化与交叉验证}

\textbf{待补章节。} 本章计划覆盖：特征合同与泄漏控制、基线（均值/岭回归）与候选模型、开发集内超参选择、按独立事件分组的交叉验证、冻结后独立留出评分、物理一致性（守恒残差）与预测精度的分别报告。历史 v3 岭回归只作为方法基线，不冒充新模型盲测。
